% 08_discussion.tex
\section{Discussion}
\label{sec:discussion}

\subsection{When to use what}

The calibration formula $S^\star(n_{\min}, \rho, \alpha)$ of
\autoref{thm:detection} reduces operational deployment to three
inputs: the minimum text length $n_{\min}$ at which marking must be
verifiable, the target false-positive rate $\alpha$, and the
watermark budget $\rho$ a deployer is willing to spend. For long-form
generated content the calibration anchors comfortably at small
$S$ with $\rho$ near $1/2$, so the marginal cost of marking is
small and detection power is high; the head-to-head comparisons in
\autoref{subsec:exp2-headline} show that MCL matches or exceeds
KGW and SWEET in this regime under matched budgets. For short
outputs $S^\star$ rises sharply with stricter $\alpha$, so deployers
targeting tight false-positive rates on short passages should
either tolerate the higher state count or shift the
specification toward a longer floor. When the expected attack
profile is dominated by paraphrase or cross-lingual rewrite
(\autoref{subsec:exp3-translation}), the universal robustness
threshold $\delta^\star$ of \autoref{thm:robustness} bounds the
worst case across all gating choices, so the practical decision
collapses to whether the deployer wants entropy gating
(KGW/SWEET-style) or the gate-free MCL specification. SWEET, which
we read as a high-entropy gate following~\citet{lee2023sweet}, is a
sensible choice when the underlying model produces sharply
peaked distributions on most positions; MCL is the more direct
choice when no such structure is assumed.

\subsection{Limitations}

\paragraph{Single budget point in headline runs.} The headline
matched-budget grid fixes $\rho = 0.5$ across all gates, models, and
domains. We have an earlier $\rho$ ablation on a smaller setup, but
re-running the sweep under the new model fleet and attack suite is
deferred. The closed-form scaling of \autoref{thm:detection} is in
$\rho$, so any deployer can re-instantiate $S^\star$ for a different
budget without re-running our experiments.

\paragraph{Self-perplexity quality metric.} Generation quality is
reported as the perplexity of each model under its own logits. We do
not report external-LM perplexity, MAUVE, or human-rater scores;
self-perplexity is a relative comparator across gates on the same
model and is conservative for absolute claims about quality.

\paragraph{Three instruction-tuned LLMs.} The headline experiments
use three open-weight instruction-tuned 7--8B parameter models. We
do not run base (non-instruction-tuned) models, smaller models, or
models above 10B. Trends across these three are consistent
(\autoref{sec:experiments}) but the population is small.

\paragraph{$k$-regular validation only at $k\in\{1,2\}$.}
\autoref{thm:k-regular} predicts $\delta^\star$ universality for
every $k$-regular topology; we validate $k=1$ (clockwork) and $k=2$
(soft-cycle) empirically. $k=3$ and beyond are deferred to the
extended version of this paper.

\paragraph{DIPPER paraphrase non-determinism.} Our paraphrase attack
uses DIPPER without an internal seed, so reruns are not bit-exact.
We mitigate by reporting per-prompt deltas from a fixed pool of
generations rather than aggregate point estimates of the attack
output, but the residual variance is real.

\subsection{Threats to validity}

The robustness analysis underlying \autoref{thm:robustness} models
the attack as i.i.d.\ uniform substitution at rate $\delta$, and
predicts post-attack $z$-scores that scale as
$(1-\delta_{\mathrm{eff}})^2$. Translation and paraphrase attacks
are not i.i.d.\ uniform substitution: they preserve syntax,
correlate substitutions across positions, and leak word identity
through morphology. The empirical fit in
\autoref{subsec:exp3-translation} is therefore best read as a
qualitative match to the theoretical scaling, with
$\delta_{\mathrm{eff}}$ an effective-edit-rate proxy rather than a
strict bound. The $(1-\delta_{\mathrm{eff}})^2$ shape predicts
ordering of attacks by severity but does not give tight
finite-sample guarantees outside the i.i.d.\ regime.
